\subsection{Vision-Language Models}

Contrastive image-text pretraining, exemplified by
CLIP~\cite{radford2021clip}, learns a shared embedding space for images and
text from large-scale web data without task-specific labels, and underlies
the image encoder used in this work's retriever
(Section~\ref{sec:method}). Building on such encoders, generative
vision-language models couple a frozen or lightly-adapted vision encoder
with a large language model to produce free-text answers conditioned on an
image. BLIP-2~\cite{li2023blip2} bridges a frozen image encoder and a
frozen language model with a lightweight Querying Transformer, substantially
reducing the number of trainable parameters relative to end-to-end
pretraining. LLaVA~\cite{liu2023llava} instead couples a vision encoder to
a language model through visual instruction tuning on
model-generated instruction-following data. This project's baseline VLM
(Section~\ref{sec:method}) is drawn from this family of pretrained,
general-domain checkpoints.

\subsection{Medical Vision-Language Models}

General-domain VLMs transfer imperfectly to medical imaging, where visual
patterns, terminology, and the stakes of an incorrect answer all differ
from natural-image domains. LLaVA-Med~\cite{li2023llavamed} adapts the
LLaVA recipe to biomedicine, continuing training on figure-caption pairs
drawn from PubMed Central and biomedical instruction-following data
generated with the assistance of a large language model, and reports
improvements over general-domain VLMs on medical visual question answering.
This line of work establishes that domain adaptation helps, but the
resulting models still answer purely from their trained parameters at
inference time --- they do not retrieve or cite supporting evidence for a
specific query, which is the gap this paper's retrieval-augmented condition
addresses.

\subsection{Retrieval-Augmented Generation}

Retrieval-augmented generation was introduced for knowledge-intensive
text tasks~\cite{lewis2020rag}, combining a pretrained sequence-to-sequence
generator with a dense vector index over a non-parametric corpus, retrieved
via maximum inner-product search. The approach's central premise ---
grounding generation in retrieved, inspectable evidence rather than relying
solely on parametric memory --- transfers naturally to a setting where
evidence is diagnostically meaningful and open to human review, as in
medical imaging. Efficient large-scale nearest-neighbor search over dense
embeddings, which the retrieval side of this pipeline depends on, is
provided by FAISS~\cite{johnson2019faiss}, used here as an in-process
CPU-based flat inner-product index (Section~\ref{sec:method}) appropriate
for a corpus of ROCOv2's scale.

\subsection{Medical VQA Benchmarks}

VQA-RAD~\cite{lau2018vqarad} is the first manually constructed medical VQA
dataset in which clinicians posed naturally occurring questions about
radiology images and supplied reference answers, spanning both closed-ended
(yes/no-style) and open-ended free-text questions across CT, MRI, and X-ray
studies. Its clinician-authored, non-synthetic questions and small,
well-defined test split make it a natural fit for the retrieval-augmented
comparison in this paper, at the cost of a test-split size small enough to
require explicit statistical significance testing
(Section~\ref{sec:experimental_setup}) before treating any accuracy delta
as reliable.

\subsection{Positioning of This Work}

Prior work has separately established that (a) domain-adapted VLMs
outperform general-domain VLMs on medical VQA~\cite{li2023llavamed}, and
(b) retrieval augmentation improves knowledge-intensive text generation in
non-medical settings~\cite{lewis2020rag}. This paper combines these two
threads directly: rather than adapting a model's weights to the medical
domain, it tests whether providing retrieved medical evidence \emph{at
inference time}, to an otherwise unmodified VLM, produces a measurable
accuracy improvement on medical VQA --- and, if so, where that improvement
does and does not appear (Section~\ref{sec:discussion}).
